\documentclass[11pt,a4paper]{article}
\usepackage[utf8]{inputenc}
\usepackage[T1]{fontenc}
\usepackage[margin=2.5cm]{geometry}
\usepackage{amsmath,amssymb}
\usepackage{booktabs}
\usepackage{array}
\usepackage{hyperref}

\newcommand{\tx}[1]{\texttt{#1}}
\newcommand{\vs}{\textvisiblespace}
\newcommand{\comment}[1]{}
\title{NSUCRYPTO 2026, Round 2\\Problem 8: ``Crypto room''}
\date{}
\author{Jacob Lu, Raymond Ouyang, Ruiyang Ye}
\begin{document}
\maketitle

\section*{Answer}

\begin{quote}
``Cryptography is typically bypassed, not penetrated.'' --- Adi Shamir
\end{quote}

\begin{itemize}
  \item \textbf{Cipher:} the 17 hexadecimal labels in the room are textbook
  RSA ciphertexts. Each one encrypts 4 ASCII bytes, read as a big-endian
  integer.
  \item \textbf{Public key:} $e=65537$ and
  \[
  N = 151\,456\,309\,706\,371 = \tx{0x89bfab89fe83} = p\,q,
  \qquad p = 5\,041\,523,\quad q = 30\,041\,777 .
  \]
  \item \textbf{Private exponent:}
  $d = e^{-1} \bmod (p-1)(q-1) = 77\,487\,966\,921\,153$.
\end{itemize}

\section{The labels}

Every label has 12 hex digits, so it is a 48-bit integer. Taken together,
the labels use every one of the 16 hex digits in positions 2 through 12. The
first digit,
however, is always one of \tx{0}--\tx{8}, and the largest label is
\tx{8649d2897de8} $\approx 1.48\cdot 10^{14}$. If the 17 labels were uniform
48-bit values, the first digit of every label would lie in $\{0,\dots,8\}$
with probability
\[
(9/16)^{17} \approx 6\cdot 10^{-5}.
\]
The labels therefore look like residues modulo a common $N$ slightly larger
than $1.48\cdot 10^{14}$. The number $65537$ is the standard RSA public
exponent, and Shamir is the ``S'' of RSA. Together these point to RSA with a
47-bit modulus. The modulus itself is not written anywhere in the picture.

\section{Recovering the modulus}

The picture shows the birth and death years of two other cryptographers on the
wall: 1523--1596 for Vigen\`ere and 1777--1855 for Gauss. These are the clue to
the modulus: its primes carry the birth years in their last four digits,
\[
p = 5\,041\,523 \equiv 1523 \pmod{10^4}, \qquad
q = 30\,041\,777 \equiv 1777 \pmod{10^4}.
\]
 With this structure, $N$ can be recovered
without guessing any plaintext.

From the previous section, $N$ lies between the largest label and
$9\cdot 16^{11}$, a window of width about $1.07\cdot 10^{13}$. Write
\[
p = 10^4 a + 1523, \qquad q = 10^4 b + 1777.
\]
Any genuine RSA modulus has primes of comparable magnitude, so we search
$p$ from $10^6$ up to $\sqrt{9\cdot 16^{11}} \approx 1.26\cdot 10^7$. This
leaves 183 prime candidates for $p$. For each of them, $q = N/p$ is forced
into the window $(\max_i c_i)/p \le q \le 9\cdot 16^{11}/p$, so $b$ ranges
over only 44\,406 values in total, of which 6\,222 make $q$ prime.

For each of the 6\,222 candidate pairs $(p,q)$ we compute
$d = e^{-1} \bmod (p-1)(q-1)$ (when the inverse exists) and decrypt one
label. Under a wrong modulus the result looks uniform in $[0,N)$, so the
chance that it is below $2^{32}$ and consists of four printable ASCII bytes
(codes 32--126) is about
\[
\frac{2^{32}}{N}\cdot\Bigl(\frac{95}{256}\Bigr)^4 \approx 5.4\cdot 10^{-7}.
\]
A second label makes a false survivor astronomically unlikely. Exactly one
pair survives:
\[
p = 5\,041\,523, \qquad q = 30\,041\,777, \qquad
N = 151\,456\,309\,706\,371.
\]
Checking all 17 labels confirms it: every label decrypts to printable bytes
and re-encrypts to itself. The whole search takes a couple of seconds.

\comment{
\paragraph*{Remark (alternative: known-plaintext recovery).}
The modulus can also be recovered without the years, at the price of
guessing a fragment of the plaintext. If $c = m^e \bmod N$, then $N$ divides
the integer $m^e - c$, so two correctly guessed plaintext--ciphertext pairs
give
\[
N \;\Big|\; \gcd\!\left(m_1^{\,e} - c_1,\; m_2^{\,e} - c_2\right),
\]
and in practice this gcd equals $N$ or a small multiple of it. Guessing that
the hidden text contains Shamir's best-known line, and trying block sizes of
3 to 6 bytes, every alignment, and all ordered pairs of labels, the pair
\[
m_1 = \tx{"ptog"},\ c_1 = \tx{0037e3e277d0}, \qquad
m_2 = \tx{"raph"},\ c_2 = \tx{5bf4286560fd}
\]
(4-byte blocks, starting 3 characters into ``Cryptography'') gives the same
$N = 151\,456\,309\,706\,371$ directly.
}
\section{Decryption}

With $d$ as above, each label $c$ decrypts to $m = c^d \bmod N$, written as
4 bytes. In the table, \vs\ marks a space.

\begin{center}
\renewcommand{\arraystretch}{1.1}
\begin{tabular}{@{}>{\ttfamily}l >{\ttfamily}l@{}}
\toprule
\normalfont Label & \normalfont Plaintext \\
\midrule
69cac604f3f6 & "Cry \\
0037e3e277d0 & ptog \\
5bf4286560fd & raph \\
7a78ae829dee & y\vs is \\
1160d891335b & \vs typ \\
25aa4a80579b & ical \\
465d1f5ae113 & ly\vs b \\
3f268417e1bd & ypas \\
751fb3515b44 & sed, \\
44a174ad7278 & \vs not \\
660301648837 & \vs pen \\
7a8d796a70ab & etra \\
5c40bf50d20b & ted. \\
3932ff60773e & "\vs -\vs \\
8649d2897de8 & Adi\vs \\
1d5a29213b82 & Sham \\
788f5bae6242 & ir\vs\vs \\
\bottomrule
\end{tabular}
\end{center}

\noindent
Joined in this order, the chunks read
\begin{center}
\tx{"Cryptography is typically bypassed, not penetrated." - Adi Shamir}
\end{center}

\end{document}
